% Exported from Math++. Compile with pdflatex.
\documentclass[11pt,a4paper]{article}
\usepackage[T1]{fontenc}
\usepackage[utf8]{inputenc}
\usepackage{lmodern}
\usepackage{amsmath,amssymb,amsthm,mathtools}
\usepackage{graphicx,booktabs,xcolor}
\usepackage[hidelinks]{hyperref}
\usepackage[margin=25mm]{geometry}
\theoremstyle{plain}
\newtheorem{theorem}{Theorem}
\newtheorem{lemma}[theorem]{Lemma}
\newtheorem{corollary}[theorem]{Corollary}
\newtheorem{proposition}[theorem]{Proposition}
\newtheorem{conjecture}[theorem]{Conjecture}
\theoremstyle{definition}
\newtheorem{definition}[theorem]{Definition}
\newtheorem{example}[theorem]{Example}
\theoremstyle{remark}
\newtheorem*{remark}{Remark}
\setlength{\parindent}{0pt}\setlength{\parskip}{0.6em}
\begin{document}
\section*{Differential equations}

Solve an equation for a function with \texttt{dsolve}.

\begin{equation*}
  \operatorname{Solve}(y' = y, y, x) = y\left(x\right) = e^{x} C
\end{equation*}
\begin{equation*}
  \operatorname{Solve}(y' = x y, y, x) = y\left(x\right) = e^{\frac{x^{2}}{2}} C
\end{equation*}
With a starting value, written as the point the solution passes through:

\begin{equation*}
  \operatorname{Solve}(y' = y, y, x, (0, 1)) = y\left(x\right) = e^{x}
\end{equation*}
A numerical solution, when there is no formula:

\begin{equation*}
  \operatorname{RK4}(x + y, x, y, 0, 1, 1, 100) \approx 3.436563656
\end{equation*}
\end{document}
